\section{Extraction from the discrete spectrum}\label{sec:extraction}

We prove a local restriction on the discrete automorphic spectrum.
All absolute values in this section refer to one fixed complex
realization.  We transport the coefficient field and the sheaf
constructions to that realization as abstract characteristic-zero linear
algebra.  The eigenvalues coming from mixed compactly supported
cohomology retain their Weil weights under this transport.

Let $G$ be split connected reductive and let $T$ be a split maximal
torus.  An \emph{ordered Bernstein weight} of a local representation
at a place $x$ is a character $t_0:X_*(T)\to\mathbf k^\times$
occurring in its Iwahori invariants,
before passage to its Weyl orbit.  Put $q_x=q^{\deg(x)}$ and define its
exponent by
\[
 \langle\lambda_0,a\rangle=\log_{q_x}|t_0(a)|.
\]
We call this exponent dominant if
$\langle\lambda_0,\alpha^\vee\rangle\ge0$ for every positive root
$\alpha$ of $G$, with the positive-translation convention of
Section~\ref{sec:categorical}.

\begin{theorem}\label{ext:discrete}
 Let $G$ be split connected reductive, let $x$ be a closed point, and
 impose Iwahori level at $x$ and arbitrary fixed compact open level
 away from $x$.  Let $\pi$ be a unitary discrete automorphic
 representation with nonzero invariants at that level.  If an ordered
 Bernstein weight $t_0$ of $\pi_x$ has dominant exponent
 $\lambda_0=\log_{q_x}|t_0|$, then
 \begin{equation}\label{ext:local-form}
 t_0\sim
 \phi\begin{pmatrix}q_x^{1/2}&0\\0&q_x^{-1/2}\end{pmatrix}c,
 \end{equation}
 where $\phi\colon\operatorname{SL}_2\to\Gd$ is a homomorphism and
 $c$ is semisimple and conjugate into a maximal compact subgroup of
 $Z_{\Gd}(\phi)$.
\end{theorem}

The proof is an induction on semisimple rank.  For a hypothetical
violating weight, a polynomial spectral filter and the rotation trace
identity force its occurrence in compact cohomology of translation
shtukas.  Integral Frobenius weights then make $\lambda_0$ rational,
allowing integral highest weights $\mu=m\lambda_0$ along its exact ray.
The Wakimoto filtration isolates the repeated highest-weight term,
and hyperspecial projections transfer its high-degree cohomology to
the spherical calculation.  A separate spherical dg model supplies
an actual chain projector onto this summand.  A finite Koszul test
preserves its highest nonzero cohomological degree, which the
Frobenius lower bound and Theorem~\ref{amp:amplitude} then bound
incompatibly as $m$ grows.  A fixed auxiliary place makes the translation
shtukas bounded even when the ray lies on a chamber wall.

\subsection{Dominant ordered weights and hyperspecial projections}

We first isolate the local algebra needed for passage from Iwahori to
hyperspecial level.  Throughout this subsection, $Q>1$ is the residue
cardinality.  Let $L=X_*(T)$ be the cocharacter lattice of a split connected
reductive group, let $W$ be its finite Weyl group, and let
\[
  \mathcal H=\mathcal H(L,Q),\qquad
  A=\mathbf k[L],\qquad Z=A^W
\]
be its Iwahori Hecke algebra, Bernstein subalgebra, and center.
We use the group-theoretic Bernstein presentation in
\cite[Lemmas~1.6.1 and~1.7.1, Remark~1.7.2, and
Equation~(1.15.2)]{HainesKottwitzPrasad}.
The coefficient field $\mathbf k$ is algebraically closed of characteristic
zero.  The unnormalized finite simple generators, denoted here by
$\mathsf T_s$, satisfy
\[
  (\mathsf T_s-Q)(\mathsf T_s+1)=0.
\]
We write $\theta_a$ for the Bernstein element corresponding to $a\in L$.
In the positive chamber of standard translations it is
\[
  \theta_a=Q^{-\ell(t^a)/2}\mathsf T_{t^a}=T_a.
\]
Thus the trivial representation has positive real Bernstein exponent.
With $\alpha$ a positive simple root of the group, the Bernstein relation
in these conventions is
\begin{equation}\label{ext:bernstein-relation}
 \theta_a\mathsf T_s-\mathsf T_s\theta_{sa}
  =(Q-1)\frac{\theta_a-\theta_{sa}}{1-\theta_{-\alpha^\vee}}.
\end{equation}
The quotient on the right belongs to $A$.

For a character $t\colon L\to\mathbf k^\times$, choose a complex
realization of its values and write
$\lambda=\log_Q|t|\in\operatorname{Hom}(L,\mathbf R)$.
We call $t$ dominant if
\begin{equation}\label{ext:dominant-character}
  \lambda(\alpha^\vee)\ge0\qquad(\alpha>0).
\end{equation}
This condition places no restriction on the central part of $\lambda$.
Put $c=Wt\in\operatorname{Spec}(Z)$ and
\[
  R=\widehat{Z}_{c},\qquad B=\mathcal H\otimes_Z R.
\]
Since $A$ is finite over $Z$, its completed scalar extension decomposes as
\[
  A\otimes_Z R\simeq\prod_{r\in Wt}\widehat A_r.
\]
Let $p_t$ be the idempotent of this product selecting the factor at $t$.
On a finite-dimensional module with central character $c$, $p_t$ selects
the generalized $A$-weight space at $t$.

Choose a base alcove in the reduced building.  Let $\mathcal V$ be its
finite set of hyperspecial vertices and let $e_v\in\mathcal H$,
$v\in\mathcal V$, be their normalized
parahoric idempotents.  In particular, at the base hyperspecial vertex,
\begin{equation}\label{ext:spherical-idempotent}
  e=\frac{\sum_{w\in W}\mathsf T_w}
          {\sum_{w\in W}Q^{\ell(w)}}.
\end{equation}
The denominator is nonzero.  All these vertices are conjugate under the
length-zero alcove stabilizer for the adjoint semisimple root datum;
the corresponding conjugate idempotents belong to $\mathcal H$ even when
those length-zero elements do not belong to its original extended Weyl
group.

\begin{lemma}\label{ext:hecke}
  Suppose that $t$ satisfies Equation~\eqref{ext:dominant-character}.
  Then, in $B$,
  \begin{equation}\label{ext:hecke-ideals}
    Bp_tB=\sum_{v\in\mathcal V}Be_vB.
  \end{equation}
  Equivalently, a simple module at the central character $Wt$ contains the
  ordered Bernstein weight $t$ if and only if at least one of the
  hyperspecial idempotents acts nontrivially on it.
\end{lemma}

\begin{proof}
  We first treat the full semisimple coweight lattice, then descend through
  a finite lattice extension, and finally pass from simple modules to
  idempotent ideals.

  \smallskip
  \noindent\emph{The full lattice.}
  Suppose first that $L$ is the direct sum of the full coweight lattice of
  the semisimple root system and a central lattice.  The latter contributes
  a central Laurent polynomial algebra and may be specialized at its
  character.  It therefore suffices to consider the semisimple factors.
  For a character $r$ put
  \[
    I(r)=\mathcal H\otimes_A\mathbf k_r.
  \]
  Its spherical line is spanned by
  $\mathbf 1_r=\sum_{w\in W}\mathsf T_w\otimes1$.
  We use the following precise principal-series theorem: for the full
  weight lattice of a root system, the spherical vector generates the
  principal-series module if and only if
  \begin{equation}\label{ext:spherical-cyclicity}
    \prod_{\alpha>0}\bigl(1-Qr(\alpha^\vee)\bigr)\ne0.
  \end{equation}
  This is \cite[Proposition~2.11(b)]{Ram}, applied to the dual root system.
  The parameter denoted by $q$ there is $Q^{1/2}$ here, and its finite
  generator is $Q^{-1/2}\mathsf T_s$.  Thus the exceptional root value in
  that proposition is $Q^{-1}$, as in
  Equation~\eqref{ext:spherical-cyclicity}.  The result includes singular
  characters.  For a product of root systems it follows by tensor product.
  Its characteristic-zero formulation over $\mathbf k$ follows by
  extension of scalars: the finite matrices determining cyclicity are
  defined over the field generated by the finitely many character values
  and $Q^{1/2}$, which admits a complex realization.

  If $r$ is dominant, every positive-coroot value has modulus at least
  one, so Equation~\eqref{ext:spherical-cyclicity} holds.  Consequently
  $I(r)$ is generated by $\mathbf 1_r$.  A simple module $S$ containing
  the generalized weight $r$ contains an ordinary eigenvector of that
  weight: a finite-dimensional module over a commutative algebra has a
  common eigenvector in each nonzero generalized weight space.  Inducing
  such a vector gives a surjection $I(r)\to S$.  The generating spherical
  vector cannot map to zero, so $eS\ne0$.

  Conversely, there is exactly one spherical simple at each central
  character.  For completeness, this follows directly from the spherical
  Satake identity $e\mathcal He=Ze$
  \cite[Sections~4.1 and~4.6, Equation~(4.6.1)]{HainesKottwitzPrasad}.
  At a central character $c'$, the
  module
  \[
    P=\mathcal He\otimes_Z\mathbf k_{c'}
  \]
  is finite-dimensional and satisfies $eP=\mathbf k e$.  Every proper
  submodule of $P$ has zero $e$-projection, since any submodule containing
  $eP$ contains the generator of $P$.  The sum of its proper submodules
  is therefore proper and is its unique maximal submodule.  Any simple
  module $S$ with $eS\ne0$ and central character $c'$ is a quotient of
  $P$: for $0\ne u\in eS$, simplicity gives $S=\mathcal H u$ and
  $eS=(e\mathcal He)u=\mathbf k u$.  This proves uniqueness.

  At $c'=Wr$, every simple quotient of $I(r)$ is spherical by the preceding
  paragraph.  Hence its unique spherical simple is such a quotient and
  contains $r$, since the inducing vector generates $I(r)$ and has a
  nonzero image in every nonzero quotient.  We have proved, for the full
  lattice and every dominant $r$,
  \begin{equation}\label{ext:full-lattice-detector}
    S[r]\ne0\quad\Longleftrightarrow\quad eS\ne0,
  \end{equation}
  where $S[r]$ denotes the generalized weight space.

  \smallskip
  \noindent\emph{A finite enlargement for a general reductive lattice.}
  Let $V=L\otimes_{\mathbf Z}\mathbf Q$.  Write
  \[
    V=V_{\mathrm{ss}}\oplus V_{\mathrm c},
  \]
  where $V_{\mathrm{ss}}$ is the rational span of the coroots and
  $V_{\mathrm c}$ is the common kernel of the roots.  Let $Q^\vee$ and
  $P^\vee$ be the coroot and full coweight lattices in $V_{\mathrm{ss}}$,
  and put
  \[
    L_{\mathrm{ss}}=L\cap V_{\mathrm{ss}},\qquad
    L_{\mathrm c}=\operatorname{pr}_{\mathrm c}(L),\qquad
    L'=P^\vee\oplus L_{\mathrm c}.
  \]
  The projection is rational, so $L_{\mathrm c}$ is a lattice.
  Root integrality gives
  $\operatorname{pr}_{\mathrm{ss}}(L)\subset P^\vee$; hence $L\subset L'$.
  The inclusion has finite index, and
  \begin{equation}\label{ext:lattice-quotient}
    \Gamma=L'/L\simeq P^\vee/L_{\mathrm{ss}}.
  \end{equation}
  Indeed the map from $P^\vee$ to $L'/L$ is surjective because
  $L\to L_{\mathrm c}$ is surjective, and its kernel is
  $L_{\mathrm{ss}}$.  In particular $Q^\vee\subset L_{\mathrm{ss}}$.

  Set $\mathcal H'=\mathcal H(L',Q)$ and $A'=\mathbf k[L']$.
  The Weyl group acts trivially on $\Gamma$, since
  $w a'-a'\in Q^\vee$ for $a'\in L'$.  The Bernstein presentation
  consequently grades $\mathcal H'$ by $\Gamma$, with degree-zero
  subalgebra $\mathcal H$.  Each component has an invertible homogeneous
  element $\theta_{a'}$, which commutes with $A$.  Thus
  \begin{equation}\label{ext:strong-grading}
    \mathcal H'=\bigoplus_{\gamma\in\Gamma}
                   \theta_{a'_{\gamma}}\mathcal H
  \end{equation}
  is a strong finite grading.  No splitting of $\Gamma$ into a subgroup
  of units is required.

  We shall also use a different set of homogeneous units.  The
  length-zero stabilizer $\Omega'_{\mathrm{ss}}$ of the semisimple base
  alcove is isomorphic to $P^\vee/Q^\vee$, and its map onto $\Gamma$ is
  surjective by Equation~\eqref{ext:lattice-quotient}.  Choose
  $\omega_\gamma\in\Omega'_{\mathrm{ss}}$ in each degree and put
  $U_\gamma=\mathsf T_{\omega_\gamma}$.  Then
  \begin{equation}\label{ext:hyperspecial-conjugates}
    e_\gamma=U_\gamma^{-1}eU_\gamma\in\mathcal H
  \end{equation}
  is the idempotent of a hyperspecial vertex.  To see the inclusion in
  $\mathcal H$, one may either use its degree zero or note that conjugation
  by $\omega_\gamma$ takes the finite Weyl subgroup at the base vertex
  to the finite parahoric Weyl subgroup at that hyperspecial vertex.
  Central translations fix the reduced building and give no further
  idempotents.  Taking all of $\mathcal V$, rather than just the vertices
  arising from these representatives, is therefore harmless.

  \smallskip
  \noindent\emph{Descent of the simple-module detector.}
  Let $S$ be a simple $\mathcal H$-module at the central character $c$.
  It is finite-dimensional, since $\mathcal H$ is finite over $Z$.
  Choose a simple quotient $S'$ of the finite-dimensional induced module
  $\mathcal H'\otimes_{\mathcal H}S$.  Adjunction gives a nonzero map
  $S\to S'$, hence an embedding on restriction.  We identify $S$ with
  its image.  Simplicity of $S'$ and the two choices of homogeneous units
  give
  \begin{equation}\label{ext:clifford-sums}
    S'=\sum_{\gamma\in\Gamma}\theta_{a'_\gamma}S
       =\sum_{\gamma\in\Gamma}U_\gamma S.
  \end{equation}
  Each summand is a simple $\mathcal H$-module, because each homogeneous
  unit normalizes the degree-zero algebra.  This also proves directly
  the needed Clifford-theoretic assertion that the restriction is a
  semisimple sum of conjugates.  More particularly,
  $\theta_{a'_\gamma}$ commutes with $A$, so all summands in the first
  sum have the same $A$-weight support.  Consequently
  \begin{equation}\label{ext:restriction-support}
    S[t]\ne0\quad\Longleftrightarrow\quad S'[t]\ne0,
  \end{equation}
  with the right-hand side interpreted as an $A$-weight space.
  The second sum gives the complementary assertion
  \begin{equation}\label{ext:restriction-spherical}
    eS'\ne0
    \quad\Longleftrightarrow\quad
    e_\gamma S\ne0\text{ for some }\gamma\in\Gamma.
  \end{equation}

  Suppose first that $S[t]\ne0$.  The nonzero $A$-weight space $S'[t]$
  contains an $A'$-weight $t'$ restricting to $t$.  Since $L$ has finite
  index in $L'$, restriction determines the real absolute-value exponent
  uniquely; thus $t'$ is dominant.  Applying
  Equation~\eqref{ext:full-lattice-detector} to $\mathcal H'$ gives
  $eS'\ne0$.  Equation~\eqref{ext:restriction-spherical} now gives a
  nonzero hyperspecial projection on $S$.

  Conversely, suppose that $e_vS\ne0$ for a hyperspecial vertex $v$.
  The idempotent $e_v$ is conjugate to $e$ by a length-zero unit of
  $\mathcal H'$, so $S'$ is spherical.  Let $Wr'$ be its $A'^W$-central
  character.  Restriction of this central character to $Z$ is $Wt$;
  hence, after replacing $r'$ by a Weyl conjugate, its restriction to
  $L$ is exactly $t$.  This representative $r'$ is dominant.  The
  full-lattice detector says that the spherical simple $S'$ contains
  $r'$, and hence contains $t$ on restriction to $A$.
  Equation~\eqref{ext:restriction-support} gives $S[t]\ne0$.  We have
  proved the asserted equivalence on all simple modules at $c$.

  \smallskip
  \noindent\emph{The completed ideals.}
  The algebra $B$ is finite over the complete local ring $R$.
  Every simple $B$-module is annihilated by the maximal ideal
  $\mathfrak m_R$: it is finite over $R$, and centrality and simplicity
  make $\mathfrak m_R S$ either zero or $S$, with the second possibility
  excluded by Nakayama's lemma.  The same assertion holds for every
  quotient of $B$.  Thus the simple $B$-modules are precisely the simple
  modules of the residue fiber to which the preceding detector applies.

  Put $J=\sum_{v\in\mathcal V}Be_vB$.  In $B/J$, the image of $p_t$
  annihilates every simple module by that detector.  It belongs to the
  Jacobson radical and is an idempotent, so it is zero.  Therefore
  $Bp_tB\subset J$.  Conversely, in $B/Bp_tB$ every image of $e_v$
  annihilates every simple module.  These images are again idempotents
  in the Jacobson radical and are zero.  This proves the reverse
  inclusion and Equation~\eqref{ext:hecke-ideals}.
\end{proof}

\begin{corollary}\label{ext:hecke-finite}
  In the notation of Lemma~\ref{ext:hecke}, let
  $N=\bigoplus_i N^i$ be a graded $B$-module whose $B$-action preserves
  degree.  No finite-generation hypothesis on $N$ is imposed.
  If $p_tN$ is finite-dimensional, then $e_vN$ is finite-dimensional for
  every $v\in\mathcal V$.  Moreover,
  \[
    p_tN^i\ne0\quad\Longrightarrow\quad
    e_vN^i\ne0\text{ for some }v\in\mathcal V.
  \]
  These assertions apply in particular to any invariant direct summand
  selected by a projector commuting with $B$.
\end{corollary}

\begin{proof}
  Equation~\eqref{ext:hecke-ideals} supplies, for each $v$, a finite
  expression $e_v=\sum_j a_jp_tb_j$ with $a_j,b_j\in B$.  Hence
  $e_vN\subset\sum_j a_j(p_tN)$, proving finite-dimensionality.
  Conversely, write $p_t$ as a finite sum of elements of the ideals
  $Be_vB$.  If all $e_vN^i$ were zero, this expression would give
  $p_tN^i=0$.  Finally, the image of a projector commuting with $B$ is a
  $B$-submodule, so the same argument applies to it.
\end{proof}

\subsection{The finite-level analytic input}

We specify the analytic facts used below.  Write $L=X_*(T)$ for a split
maximal torus and put
\[
 L^0=\{a\in L:\langle\chi,a\rangle=0
                    \text{ for every }\chi\in X^*(G)\}.
\]
This is a saturated sublattice spanning the semisimple directions.
The annihilator of $L^0$ in the dual torus is $Z(\Gd)^\circ$.
Consequently two ordered parameters have the same character on $L^0$
if and only if they differ by an element of $Z(\Gd)^\circ$.

Choose a free group $\Xi$ of rational central modifications whose degree
image is cocompact in the free central degree lattice and intersects
degree zero trivially, as in Section~\ref{sec:rotation}.
For a compact open level $K$, let
\[
 \mathcal L_K=L^2\bigl(G(F)\backslash G(\mathbb A_F)/(K\Xi)\bigr).
\]
The measure on its point basis includes the usual automorphism factors.
An arbitrary unitary automorphic central character can be treated in
this model after a unitary unramified twist: a character on the lattice
$\Xi$ extends to the degree group of $G$, in which its image has finite
index, since the circle group is divisible.  Such a twist has zero real exponent and does not affect the
assertion we shall prove.

\begin{proposition}[Finite-level spectral facts]\label{ext:spectral}
 At a fixed compact open level the following assertions hold.
 \begin{enumerate}[label=\textup{(\roman*)}]
 \item The discrete summand $\mathcal D_K$ of $\mathcal L_K$ is
 finite-dimensional.  Its orthogonal complement admits a Hecke-equivariant unitary
 realization as a closed invariant subspace of a finite sum of
 direct integrals of normalized parabolic inductions
 \[
    \operatorname{Ind}_{P(\mathbb A_F)}^{G(\mathbb A_F)}
                 (\sigma\otimes\chi),
 \]
 where $P$ is proper, $\sigma$ belongs to the discrete spectrum of its
 Levi $M$, and $\chi$ ranges over a compact torus of unitary unramified
 characters, subject to the central condition imposed by $\Xi$.
 All fixed-level fiber dimensions are bounded.
 \item Let $Q_b$ be the orthogonal projection onto the point basis in an
 HN ball of radius $b$.  There are constants $C,c>0$ such that
 \begin{equation}\label{ext:discrete-tail}
   \|(1-Q_b)v\|\le C e^{-cb}\|v\|
       \qquad(v\in\mathcal D_K, b\ge0).
 \end{equation}
 \item Suppose that the level at $x$ is Iwahori.  The ordered Bernstein
 weights of a fiber induced from $M$ have the form
 \begin{equation}\label{ext:induced-weights}
      w(t_M z),\qquad w\Phi_M^+\subset\Phi_G^+,
 \end{equation}
 where $t_M$ is an ordered Bernstein weight of $\sigma_x$ and $z$ is the
 local value of $\chi$.  There are finitely many such weight families,
 and their generalized-weight lengths are uniformly bounded.
 \end{enumerate}
\end{proposition}

\begin{proof}
 We use the following precise form of Langlands' spectral theorem.
 At a fixed compact open level, the discrete pairs $(M,\sigma)$ with
 nonzero fixed-level induced vectors form finitely many classes under
 Weyl association and unramified twisting.  Their compact induced
 models are finite-dimensional, their unitary twisting spaces are
 compact real tori modulo finite stabilizers, and Eisenstein wave
 packets on these spaces give the Hecke-equivariant Hilbert spectral
 decomposition.  This formulation is stated in
 \cite[Theorem III.7 and the preceding definitions, pp.~33--36]{LLafforgueSpectral},
 where it is identified with \cite[Section VI.2.1]{MW}; the
 function-field treatment in the English edition is
 \cite[Section VI.2.7]{MW}.  The statement includes discrete inducing
 representations.  If the given level is not contained in the chosen
 standard maximal compact subgroup, intersect it with that subgroup
 first and then take invariants under the original compact level.
 This replaces the full compact models by closed invariant subspaces
 and preserves all the asserted finiteness and norm bounds.

 We explain the effect of $\Xi$.  Unramified characters are characters
 of a free degree lattice.  On the central degree lattice of $G$, the
 condition that $\sigma\otimes\chi$ be trivial on $\Xi$ fixes finitely
 many characters because the image of $\Xi$ has finite index.
 For $M=G$ this leaves finitely many points of each twisting torus.
 There are finitely many pairs by the quoted theorem and their
 fixed-level models have finite dimension, so $\mathcal D_K$ is finite.
 For a proper Levi $M$, the condition cuts its unitary parameter torus
 into finitely many translates of a closed subtorus, hence leaves a
 compact parameter space.  The wave-packet construction allows finite
 Weyl identifications between parameters; equivalently it realizes
 the corresponding Hilbert space inside finitely many full induced
 compact models.  This proves (i) in precisely the form needed here.

 Boundedness of fixed-level induction ranks can also be read directly
 from the compact picture.  The quotient of a maximal compact subgroup
 by its intersection with the inducing parabolic and the prescribed
 open level is finite.  Its finitely many representatives replace the
 level on $\sigma$ by finitely many compact open subgroups of $M$.
 Their invariant spaces are finite-dimensional.  Unitary unramified
 twists are trivial on these compact subgroups, so the same dimension
 bounds hold throughout each torus.

 For (ii), use the exact constant-term assertion on deep
 function-field Siegel sectors
 \cite[Chapter I, Section 5.9, p.~107]{MorrisI}; see also
 \cite[Section I.2.7]{MW}.  Once the root coordinates outside a Levi
 are sufficiently large, a function at the prescribed compact level
 equals its constant term for that parabolic.  The threshold is
 uniform at a fixed level and on the bounded transverse data.
 The constant-term exponent theorem says that translating an
 automorphic form's normalized constant term in the split-center
 degree variables gives a finite sum of characters times polynomials
 \cite[Section 6.2, pp.~200--201, Equations (6.1)--(6.3)]{BernsteinLapid};
 see also \cite[Sections I.3.1--I.3.2 and I.4.2]{MW}.
 We need the strict decay consequence, and give its elementary proof
 to include all boundary directions.

 Partition a Siegel domain according to which simple-root coordinates
 are above the constant-term threshold.  The remaining roots define
 the Levi; their bounded coordinates range over bounded data modulo
 its center.  After passing to finite-index degree lattices, finitely
 many transverse values and translates of positive orthants suffice.
 Take the image of each sufficiently deep sector in
 $P(F)\backslash G(\mathbb A)/(K\Xi)$, with its quotient Iwasawa
 measure $\delta_P(m)^{-1}\,du\,dm\,dk$.
 The compact unipotent quotient $U(F)\backslash U(\mathbb A)$ has
 volume one, and the function equals its constant term $f_P$ on
 this sector.  The deep-sector separation statement
 \cite[Chapter I, Section 4.2, p.~101]{MorrisI} puts every additional
 rational identification within a sector in $P(F)$, which has
 already been quotiented out.  The finitely many bounded transverse
 charts of \cite[Chapter I, Section 4.6, pp.~102--103]{MorrisI}
 therefore give bounded residual multiplicity for the map to the
 automorphic quotient.  The sector's weighted squared norm is
 consequently at most a fixed multiple of the global squared norm.
 This statement uses the parabolic quotient and its measure;
 the integral-unipotent multiplicity in a raw subset of
 $G(\mathbb A)$ need not be bounded.

 Absorb the finitely many transverse measure factors into a norm
 on a finite-dimensional space $V$.  Multiplication by
 $\delta_P^{-1/2}$ then normalizes the constant term to an
 exponential-polynomial sequence
 \[
      F\colon\mathbf N^r\longrightarrow V,
      \qquad \dim V<\infty,
 \]
 and the preceding norm comparison gives
 $F\in\ell^2(\mathbf N^r,V)$.

 Let $E$ be the span of all forward coordinate shifts of $F$.
 It is finite-dimensional because $F$ is exponential-polynomial.
 Each coordinate shift $S_j$ preserves $E$ and is a contraction for
 its inherited $\ell^2$ norm.  It has no eigenvalue of modulus one:
 if $S_jh=zh$, $|z|=1$, and $h\ne0$, a nonzero slice in the other
 coordinates gives a coordinate ray of constant nonzero norm,
 contradicting square summability.  Thus choose $\rho<1$ greater
 than the moduli of every eigenvalue of every $S_j|_E$.
 Finite-dimensional Jordan estimates and commutation of the shifts
 give
 \[
   \|S_1^{n_1}\cdots S_r^{n_r}\|\le C\rho^{n_1+\cdots+n_r}.
 \]
 Evaluation at the origin is bounded in the $\ell^2$ norm, so this
 proves exponential decay of $F$.  Summing over the polynomially
 many lattice points of a given radius gives an exponential squared
 norm tail.  Bounded sectors contribute none for large radius.
 The reduction height and the HN radius are comparable up to fixed
 constants, by the degree description of the canonical reductions.
 This proves Equation~\eqref{ext:discrete-tail} for each discrete form.
 A basis of the finite-dimensional space $\mathcal D_K$ makes the
 constants uniform on that space.

 Finally, compatibility of normalized parabolic induction with Iwahori
 invariants identifies the Iwahori invariants of the local induced representation with the
 corresponding induced Hecke module
 \cite[Section~5.8, Theorem~8]{PrasadBernstein}.
 We use left modules and the positive-translation convention of
 Equation~\eqref{ext:bernstein-relation}; the conversion from the
 source's right-module convention is described in its Section~7.
 The Bernstein presentation now proves (iii). The ambient Hecke
 algebra is free as a right module over the Levi Hecke algebra,
 with basis indexed by the representatives $w$ for which
 $w\Phi_M^+\subset\Phi_G^+$.  Commuting a Bernstein element past these
 basis elements by Equation~\eqref{ext:bernstein-relation} is triangular
 in Bruhat order, and its diagonal characters are exactly
 Equation~\eqref{ext:induced-weights}.  The real exponent of a unitary
 $\sigma$ is zero in the split central directions of $M$.  The finite
 list of inducing data and the bounded fiber dimensions from (i)
 give the last assertions of (iii).
\end{proof}

We shall use the following elementary uniform bound to handle the fact
that Bernstein translations need not be normal operators.

\begin{lemma}\label{ext:matrix-bound}
 Let $B,r_0>0$ and $R_0\in\mathbf N$ be fixed.  For every $r>r_0$ there
 is a constant $C$ with the following property: if $A$ is an operator
 on a Hilbert space of dimension at most $R_0$, with
 $\|A\|\le B$ and all eigenvalues of modulus at most $r_0$, then
 \[
       \|A^n\|\le Cr^n\qquad(n\ge0).
 \]
\end{lemma}

\begin{proof}
 On the circle $|z|=r$, the determinant of $z-A$ has absolute value at
 least $(r-r_0)^{\dim A}$.  In an orthonormal basis, its adjugate has a
 bound depending only on $B,r,R_0$.  The resolvent is therefore uniformly
 bounded on that circle.  The contour formula for $A^n$ gives the stated
 estimate, after increasing $C$ to cover the zero-dimensional case.
\end{proof}

\subsection{A dominant weight forces high cohomology}

To set up the induction for Theorem~\ref{ext:discrete}, suppose that it
is known in smaller semisimple rank and that $t_0$ violates its conclusion.
We work at a level, and after a unitary central twist, for which $\pi$ is contained in
$\mathcal D_K$.  Its central exponents are zero.  In particular
$\lambda_0$ lies in the semisimple directions, and it is nonzero: an
exponent-zero semisimple parameter has the form in
Equation~\eqref{ext:local-form} with the trivial homomorphism $\phi$.

Choose a second closed point $y\ne x$ where $\pi$ is unramified, refine
its level to an Iwahori, and fix an ordered Bernstein weight $t_{0,y}$
there. Shrink the level away from $x,y$ to a product of principal
congruence subgroups.  Every compact open subgroup contains such a
product, so the resulting level
$K=I_x\times I_y\times K^{x,y}$ still has $\pi^K\ne0$.
The factor $K^{x,y}$ is full level along an effective divisor disjoint
from $x,y$, as required by the geometric constructions.
All these levels now remain fixed.

We give the finite-level tensor argument ensuring simultaneous
occurrence of the two chosen local weights; the general tensor-product
theory is due to Flath \cite{FlathTensorProducts}.
The finite-dimensional space $V=\pi^K$ is simple over the product
Hecke algebra $\mathcal H_x\otimes\mathcal H_y\otimes\mathcal H^{x,y}$.
Indeed any nonzero vector generates $\pi$ under the adelic group by
irreducibility; averaging its translates over $K$ shows that it
generates $V$ under this Hecke algebra.
Let $A_x,A_y$ be the images of the two local Hecke algebras in
$\operatorname{End}(V)$. The subspace
$\operatorname{rad}(A_x)V$ is stable under the entire product algebra,
so is either zero or $V$. The second possibility contradicts
nilpotence of the radical; hence $A_x$ is semisimple.
Every element of its center commutes with the product algebra and
therefore acts as a scalar on $V$. Thus $A_x$ is a single matrix
algebra. The same argument applies to $A_y$.
The commuting matrix-algebra actions consequently give
\[
                  V\simeq U_x\otimes U_y\otimes W,
\]
where $U_x,U_y$ are their simple modules and $W\ne0$ is the
multiplicity space. The simple local types are independent of the
auxiliary level: refinement embeds the finite invariant spaces,
whose isotypic restrictions therefore have the same simple type.
Taking the union over auxiliary levels identifies these types with
the local Iwahori modules $\pi_x^{I_x}$ and $\pi_y^{I_y}$.
The tensor of their nonzero generalized weight spaces at $t_0$ and
$t_{0,y}$ is therefore nonzero. Thus $(t_0,t_{0,y})$ occurs on the
finite-dimensional discrete space at this refined level.
In the following, weights are compared on $(L^0)^2$ unless explicitly
stated otherwise.

\begin{lemma}\label{ext:spectral-gap}
 There is a Laurent polynomial $P_0$ in the degree-zero Bernstein
 translations at $x$, nonzero at $t_0$, and an open cone of dominant
 regular directions arbitrarily close to $\lambda_0$, such that the
 following holds.  For every fixed auxiliary label $\zeta$ at $y$
 and every sufficiently large label $\mu$ in this cone, every continuous
 joint weight surviving $P_0$ has strictly smaller modulus on
 $T_{(\mu,\zeta)}$ than the target joint weight.
\end{lemma}

\begin{proof}
 Complexify the twisting torus of each family in
 Proposition~\ref{ext:spectral}.  Its ordered $x$-weights form a coset
 of an algebraic subtorus.  There are finitely many such cosets.
 A coset not containing the target on $L^0$ can be annihilated by a
 Laurent polynomial nonzero at the target.  Raise that polynomial to
 a power exceeding the uniformly bounded generalized-weight length,
 and multiply over these cosets.  The resulting $P_0$ kills their
 actions, including generalized weights.

 Consider a remaining family, induced from $M$ and shuffled by $w$.
 Let $C_M$ be its real twisting-center subspace after that shuffle, and
 let $\nu$ be the real exponent of its unitary contour.  Unitarity of
 the inducing central character gives $\nu\perp C_M$.  Membership of
 the target in its complexified coset gives
 $\lambda_0-\nu\in C_M$.  The possible ambiguity from working on $L^0$
 is a connected-central factor of $G$ and is already contained in the
 twisting center of $M$.  Therefore
 \begin{equation}\label{ext:orthogonal-projection}
  \nu=\operatorname{pr}_{C_M^\perp}(\lambda_0),\qquad
  (\lambda_0,\nu)=\|\lambda_0\|^2
                       -\|\operatorname{pr}_{C_M}(\lambda_0)\|^2.
 \end{equation}
 If the second subtracted term were zero, the twisting factor taking
 the inducing parameter to $t_0$ would be unitary.  Moreover,
 $w\Phi_M^+\subset\Phi_G^+$ and dominance of $\lambda_0$ would make
 the corresponding inducing weight dominant for $M$.  The induction
 hypothesis would give Equation~\eqref{ext:local-form} for that weight,
 and hence for its unitary central twist and Weyl conjugate $t_0$.
 This contradicts the choice of $t_0$.  Thus every remaining family
 has
 \begin{equation}\label{ext:strict-projection}
       (\lambda_0,\nu)<\|\lambda_0\|^2.
 \end{equation}
 The finite list of strict inequalities persists in a small open cone
 of directions adjacent to $\lambda_0$ and inside the regular dominant
 chamber.  The $y$-exponents run through a finite list on the unitary
 contours.  Their pairings with the fixed $\zeta$ are bounded constants,
 whereas the strict $x$-gap grows linearly with $\mu$.  This proves the
 assertion for all sufficiently large $\mu$ in a slightly smaller cone.
\end{proof}

Choose once and for all a sufficiently divisible regular positive label
$\zeta\in L^0$ such that all the finitely many parabolic pairings
involving $\zeta$ alone are nonzero.  This is possible by avoiding a
finite collection of rational hyperplanes.  The divisibility includes
the fixed multiples required in the rotation calculation and the even
multiple needed to remove convention signs.  We impose the same fixed
divisibility on the varying $\mu$.

\begin{proposition}\label{ext:lower}
 Under the induction assumption and the choice of a violating weight
 above, let $u=(\mu,\zeta)$, where $\mu$ is sufficiently large in the
 cone of Lemma~\ref{ext:spectral-gap} and $u$ satisfies
 Equation~\eqref{rot:genericity}.  Then $M(u)$ contains the joint
 Bernstein character $(t_0,t_{0,y})$ on $(L^0)^2$.  At some full ordered
 extension $(t_x,t_y)$ of that character it has nonzero cohomology in a
 degree $i$ satisfying
 \begin{equation}\label{ext:degree-lower}
    i\ge 2d_x\langle\lambda_0,\mu\rangle
                       +2\log_q|\zeta(t_{0,y})|.
 \end{equation}
\end{proposition}

\begin{proof}
 The complex $M(u)$ has finite-dimensional total cohomology by
 Lemma~\ref{rot:bounded}.  If the target character did not occur,
 a Laurent polynomial in the two degree-zero Bernstein algebras would
 annihilate its entire cohomology and remain nonzero at the target.
 Multiply it by $P_0$ and by polynomial factors killing every other
 discrete joint character.  Powers larger than the finite generalized
 weight lengths kill the corresponding generalized spaces.  Denote
 the resulting polynomial by $P$.  It is nonzero at the target, kills
 $M(u)$, and on the discrete spectrum retains only the target joint
 generalized weight space $\mathcal D_0$.

 Multiply $P$ by a Bernstein monomial with sufficiently large positive
 labels at both places.  This changes none of these properties, since
 all Bernstein monomials are invertible.  We can consequently write
 \[
       P=\sum_{v}c_vT_v
 \]
 with a finite set of regular positive degree-zero labels $v$.  For
 the remainder of this argument $u$ and $P$ are fixed and only the
 positive integer $n$ varies.  Apply Theorem~\ref{rot:trace} to each
 term.  The normalized version of that identity gives, for every
 sufficiently large $n$,
 \begin{equation}\label{ext:filtered-diagonal}
   \sum_b(T_u^nP)_{b,b}
       =\operatorname{Tr}_{\mathrm{alt}}
                    (\mathrm{Fr}^nP\mid M(u))=0.
 \end{equation}
 Here $b$ indexes the normalized point basis of $\mathcal L_K$;
 normalization by the square root of its measure does not change
 diagonal matrix entries.  The translation $u$ is even; each insertion $T_v$ is in the
 function normalization of that theorem, including its known parity
 correction to the perverse trace.  No evenness condition is imposed
 on $v$.  Only
 this ordinary diagonal sum is used; the operator is not asserted
 to be trace class on all of $\mathcal L_K$.

 The support bound in Theorem~\ref{rot:trace} puts every contributing
 diagonal entry in an HN ball of radius $Bn$ for a fixed $B$.  Increase
 $B$ if necessary and write $Q_n=Q_{Bn}$.  The left side of
 Equation~\eqref{ext:filtered-diagonal} is therefore
 $\operatorname{Tr}(Q_nT_u^nP Q_n)$.  By
 Proposition~\ref{amp:geometry}, $\operatorname{rank}Q_n$ is polynomial
 in $n$.

 Put
 \[
     a=\mu(t_0)\zeta(t_{0,y}),\qquad R=|a|>0.
 \]
 Lemma~\ref{ext:spectral-gap} gives an $r_0<R$ bounding the moduli of
 all continuous eigenvalues of $T_u$ on the image of $P$.  Fixed Hecke
 operators are uniformly bounded on all unitary fibers by their
 $L^1$ norms.  The image of $P$ in each fiber is invariant under $T_u$,
 has bounded dimension, and inherits its Hilbert norm.  Choose
 $r_0<r<R$ and apply Lemma~\ref{ext:matrix-bound} on this image.
 We obtain
 \[
       \|T_u^nP\|_{\mathrm{cont}}\le C r^n.
 \]
 This bound is uniform even where the dimension of the image changes.
 Hence the continuous contribution to the compressed trace is at
 most $\operatorname{rank}(Q_n)Cr^n=o(R^n)$.

 On $\mathcal D_0$, simultaneous triangularization of the commuting
 Bernstein operators gives
 \begin{equation}\label{ext:discrete-leading-trace}
    \operatorname{Tr}_{\mathcal D_K}(T_u^nP)
        =\dim(\mathcal D_0)P(t_0,t_{0,y})a^n.
 \end{equation}
 The coefficient is nonzero; generalized-weight nilpotents do not
 change the diagonal of these matrices.  The norm of $T_u^nP$ on this
 finite-dimensional space is at most $C(1+n)^N R^n$ for fixed $C,N$.
 Equation~\eqref{ext:discrete-tail} therefore makes the difference
 between its full trace and its $Q_n$-compressed trace $o(R^n)$.
 Decomposing the Hilbert space into its orthogonal discrete and
 continuous summands, Equations~\eqref{ext:filtered-diagonal} and
 \eqref{ext:discrete-leading-trace} would imply
 \[
      0=\dim(\mathcal D_0)P(t_0,t_{0,y})a^n+o(R^n),
 \]
 an impossibility.  Thus the target character occurs in $M(u)$.

 Since the full Bernstein torus also acts on this finite-dimensional
 space, the target character on $(L^0)^2$ has at least one full ordered
 extension $(t_x,t_y)$ in its support.  The cyclic-slide identity of
 Proposition~\ref{cat:cyclic} identifies normalized Frobenius on that
 generalized weight space with the translation $T_u$, up to the
 already harmless convention sign.  Every one of its Frobenius
 eigenvalues has modulus
 \[
   |\mu(t_x)\zeta(t_y)|
       =q^{d_x\langle\lambda_0,\mu\rangle
                                  +\log_q|\zeta(t_{0,y})|}.
 \]
 Recall the weighted translation length
 $L_u=d_x\ell(t^\mu)+d_y\ell(t^\zeta)$ and the normalization in
 Equation~\eqref{rot:complex}.  The reduction in the
 proof of Lemma~\ref{rot:bounded} applies the compact-support weight
 theorem \cite[Corollary~3.3.4]{DeligneII} to $Y_u$: the weights of
 the constant sheaf's $H_c^j$ are integers at most $j$.  Since
 \[
    H^i(M(u))=H_c^{i+L_u}(Y_u,\mathbf k)(L_u/2),
 \]
 the weights of this normalized group are integers at most $i$.
 Comparing with the displayed Frobenius modulus proves
 Equation~\eqref{ext:degree-lower}.
\end{proof}

\begin{corollary}\label{ext:rationality}
 The semisimple exponent $\lambda_0$ of a violating weight in this
 induction argument is rational.
\end{corollary}

\begin{proof}
 The Frobenius weights of $M(u)$ are integers by the weight theorem
 and normalization just used in Proposition~\ref{ext:lower}.  Set $b=2\log_q|\zeta(t_{0,y})|$.  The occurrence in
 Proposition~\ref{ext:lower}, not merely its degree inequality, gives
 \begin{equation}\label{ext:integral-pairings}
       2d_x\langle\lambda_0,\mu\rangle+b\in\mathbf Z
 \end{equation}
 for every sufficiently large allowed label in the open cone.
 Let $L_1\subset L^0$ be the fixed finite-index lattice imposing the
 divisibility conditions.  Given $a\in L_1$, choose $\mu$ sufficiently
 deep in that cone that both $\mu$ and $\mu+a$ belong to it, and avoid
 the finitely many affine hyperplanes excluded by
 Equation~\eqref{rot:genericity} for both labels.  Such labels exist:
 a full-dimensional cone has lattice points arbitrarily far from any
 prescribed finite collection of proper affine hyperplanes.
 Subtracting the two instances of
 Equation~\eqref{ext:integral-pairings} gives
 $2d_x\langle\lambda_0,a\rangle\in\mathbf Z$.  These pairings on the
 full-rank lattice $L_1$ prove the assertion.
\end{proof}

\subsection{The extreme summand and specialization}

Use the rational metric to identify the exponent direction with a label
direction.  By Corollary~\ref{ext:rationality} we can now take
\begin{equation}\label{ext:exact-ray}
                 \mu=m\lambda_0
\end{equation}
for arbitrarily large divisible integers $m$.  These labels may lie on a
wall.  The spectral gap proved above also holds on this exact ray by
Equation~\eqref{ext:strict-projection}.  The fixed choice of $\zeta$
makes Equation~\eqref{rot:genericity} hold for every sufficiently large
such $m$: if a main-ray parabolic pairing vanishes, its auxiliary pairing
does not; if it does not vanish, the combined pairing has at most one
exceptional value of $m$.  Thus Proposition~\ref{ext:lower} and its proof
apply along this ray as well.

\begin{lemma}\label{ext:unique-extreme}
 For $\mu$ as in Equation~\eqref{ext:exact-ray}, the only weight $a$ of
 the irreducible representation $V_\mu$ maximizing
 $\langle\lambda_0,a\rangle$ is $a=\mu$, and its multiplicity is one.
 Consequently the only maximizing tuple in $V_\mu^{\otimes d_x}$ is
 $(\mu,\ldots,\mu)$.
\end{lemma}

\begin{proof}
 Every weight has the form $\mu-\sum_\alpha n_\alpha\alpha$ with
 $n_\alpha\ge0$ in the positive simple roots of the dual group.
 Equality of its pairing with that of $\mu$ forces every root appearing
 in this difference to be orthogonal to $\lambda_0$.  The highest
 weight $\mu$ has zero pairing with the coroots of that Levi root
 system.  Its highest line therefore generates the trivial
 representation of the corresponding derived Levi.  Applying only
 its lowering operators produces no different weight.  The maximizing
 face consists of the highest line alone.  Additivity of the pairing
 proves the assertion about tuples.
\end{proof}

Let $(t_x,t_y)$ be one of the full extensions supplied by
Proposition~\ref{ext:lower} for the chosen $m$.  It can vary with $m$,
but $t_x/t_0\in Z(\Gd)^\circ$.  Put $d=d_x$ and consider the Iwahori
trace complex
\[
   C_\mu=\mathsf K\bigl(Z(V_\mu)^{\boxtimes d};J_\zeta\bigr),
\]
where the semicolon records the fixed auxiliary label at $y$ and all
other fixed levels and kernels.  Let $\mathcal H_x$ be the affine
Hecke algebra at $x$ with coefficients in $\mathbf k$, let $Z_x$ be
its center, and set
\[
  R_x=\widehat{(Z_x)}_{Wt_x},\qquad
  \widehat H_\mu=H^*(C_\mu)\otimes_{Z_x}R_x.
\]
This is flat scalar extension of a graded cohomology module.
No dg $Z_x$-module structure on the Iwahori complex is used.
The auxiliary label at $y$ remains fixed; no projection onto a
$y$-weight is needed in this step.

We use a different operator here from the full normalized Frobenius
in Proposition~\ref{ext:lower}.  On the repeated highest $x$-tuple,
the $d$-fold $x$-label slide has character $d\mu$, whereas full
Frobenius on $M(\mu,\zeta)$ has joint character $(\mu,\zeta)$.
The selected-place slide separates Wakimoto terms; the lower degree
bound retains the single residue-degree factor $d_x$ in
Equation~\eqref{ext:degree-lower}.

\begin{proposition}\label{ext:extreme-transfer}
 Fix a positive even integer $\kappa$ removing the slide convention
 signs, and let $S_x$ be the $\kappa$th power of the $d$-fold
 $x$-label slide.  There is a polynomial projector
 $\varepsilon_{\mathrm{ext}}$ on $\widehat H_\mu$, commuting with
 $\mathcal H_x\otimes_{Z_x}R_x$, such that every
 \[
        N_v=e_v\varepsilon_{\mathrm{ext}}\widehat H_\mu
 \]
 has finite-dimensional total graded space.  For some hyperspecial
 vertex $v$, it is nonzero in a degree satisfying
 Equation~\eqref{ext:degree-lower}.  The transfer identifies $N_v$
 with the extreme $S_x$-part of
 $H^*(C_{\mu,v}^{\mathrm{sph}})\otimes_{Z_x}R_x$, where
 $C_{\mu,v}^{\mathrm{sph}}
   =\mathsf K_v(V_\mu^{\boxtimes d};J_\zeta)$.
\end{proposition}

\begin{proof}
 Use the finite Wakimoto filtration of
 Theorem~\ref{cat:central} and its cohomology spectral sequence.
 Its first page has labels $\mathbf a=(a_1,\ldots,a_d)$ of weights
 of $V_\mu$, with their weight-space multiplicities.  The Bernstein
 algebra and $S_x$ act on this spectral sequence.  The commutation
 statements of Propositions~\ref{cat:cyclic} and \ref{cat:hecke}
 also make $S_x$ commute with the full Hecke action on its abutment.
 Tensoring the pages and the abutment with the flat algebra $R_x$
 gives a convergent spectral sequence with its finite filtration.
 Write $\Theta_b$ for the Bernstein element of a label $b$.
 Taking the $d$-fold slide first returns each tuple $\mathbf a$ to
 its original ordering.  For $b=a_1+\cdots+a_d$, the Wakimoto
 interchange formula of Theorem~\ref{cat:central} identifies $S_x$
 on this graded term with
 \[
                  \Theta_b^\kappa\otimes U_{\mathbf a},
 \]
 where $U_{\mathbf a}$ acts on the finite tensor product of
 weight-multiplicity spaces.  The chosen power removes the parity
 signs, and the trivial semisimplified Weil structures make
 $U_{\mathbf a}$ unipotent.  The two factors commute.
 After specializing the Bernstein algebra at an ordering $r$,
 its only possible eigenvalue is
 \begin{equation}\label{ext:graded-slide}
                     \left(\prod_{j=1}^{d}a_j(r)\right)^\kappa.
 \end{equation}

 Here is an annihilator construction entirely on cohomology.
 The completed Bernstein algebra is finite over $R_x$.
 For each of the finitely many labels $\mathbf a$, put
 \[
      D_{\mathbf a}(z)=\prod_{w\in W}
           \bigl(z-\Theta_{wb}^\kappa\bigr).
 \]
 Its coefficients are invariant, and hence lie in $R_x$.
 One factor gives $D_{\mathbf a}(\Theta_b^\kappa)=0$.
 Therefore $D_{\mathbf a}(S_x)$ lies in the ideal generated by
 $1\otimes(U_{\mathbf a}-1)$ in the commutative algebra of these
 operators.  Choose $n_{\mathbf a}$ with
 $(U_{\mathbf a}-1)^{n_{\mathbf a}}=0$.  The monic polynomial
 $Q_{\mathbf a}=D_{\mathbf a}^{n_{\mathbf a}}\in R_x[z]$
 annihilates the entire first-page term, regardless of the size
 of its cohomology.  Its residual roots are exactly the values in
 Equation~\eqref{ext:graded-slide} at the orderings of $t_x$.
 Put $Q_0=\prod_{\mathbf a}Q_{\mathbf a}$.  This annihilates the
 entire first page, hence every later page and the associated graded
 of the abutment.  If the filtration has $b_\mu$ steps, then
 $Q=Q_0^{b_\mu}$ annihilates $\widehat H_\mu$: each application of
 $Q_0(S_x)$ lowers its filtration by at least one step.  The same
 $Q$ also annihilates every page.

 The maximal modulus of these residual roots is independent of
 the ordering of $t_x$, by Weyl invariance of the weight polytope.
 Hensel factorization over $R_x$ gives
 $Q=Q_+Q_-$, with $Q_+$ collecting all roots of that maximal modulus
 and $Q_-$ all remaining roots.  The factors are relatively prime.
 Their Bezout identity defines an idempotent
 $\varepsilon_{\mathrm{ext}}$, polynomial in $S_x$, on every page,
 the filtered cohomology, and its abutment.  On the abutment it
 commutes with the full Hecke algebra, since $S_x$ does and its
 polynomial coefficients are central.

 Apply also the ordered idempotent $p_{t_x}$ from
 Lemma~\ref{ext:hecke}.  It commutes with the spectral-sequence
 maps and with $\varepsilon_{\mathrm{ext}}$.
 Lemma~\ref{ext:unique-extreme} shows that exactly one label tuple
 survives on the projected first page: every entry is $\mu$.
 Its multiplicity is one.  The Weil normalization and interchange
 compatibility in Theorem~\ref{cat:central} identify this term with
 the $t_x$-ordered part of $H^*(M(\mu,\zeta))$.
 Central factors cause no additional eigenvalues, since
 $\mu\in L^0$ and $V_\mu$ is trivial on $Z(\Gd)^\circ$.
 There is no other surviving filtration step to support a
 differential.  It follows that
 $p_{t_x}\varepsilon_{\mathrm{ext}}\widehat H_\mu$ is
 finite-dimensional and contains the nonzero high-degree piece
 of Proposition~\ref{ext:lower}.  Flat completion preserves this
 finite-dimensional generalized central-character piece.

 Apply Corollary~\ref{ext:hecke-finite} to the graded module
 $\varepsilon_{\mathrm{ext}}\widehat H_\mu$.
 Each $N_v$ is finite-dimensional, and some $N_v$ retains a nonzero
 piece in the same high degree.  The cohomology transfer in
 Proposition~\ref{cat:hecke} identifies
 $e_v\widehat H_\mu$ with
 $H^*(C_{\mu,v}^{\mathrm{sph}})\otimes_{Z_x}R_x$ and intertwines
 $S_x$ and the center.  This proves the stated identification.
 The finite set of hyperspecial variants is among those allowed
 in Theorem~\ref{amp:amplitude}.
\end{proof}

The next step supplies the actual complex needed for specialization.
It uses the spherical dg model separately from the Iwahori
cohomology construction.

\begin{proposition}\label{ext:spherical-projector}
 For each $v$, the genuine spherical $Z_x$-dg module
 $C_{\mu,v}^{\mathrm{sph}}$ has, after flat scalar extension to
 $R_x$, a dg direct summand $E_{\mu,v}$ with
 $H^*(E_{\mu,v})=N_v$.  Its chain projector is polynomial in the
 spherical slide $S_x$ and commutes with finite Koszul tests over
 $R_x$.
\end{proposition}

\begin{proof}
 Use Proposition~\ref{cat:loops}.  On the finite representation
 bundle $V_\mu^{\otimes d}$, the $d$-fold slide is the degree-zero
 closed matrix
 \[
       \rho_\mu(s)=V_\mu(s)^{\otimes d}
 \]
 over $\mathcal O(\Gd)$.  Therefore
 \[
       P_\mu(z)=\det\bigl(z-\rho_\mu(s)^\kappa\bigr)
                  \in\mathcal O(\Gd)^{\Gd}[z]=Z_x[z]
 \]
 satisfies $P_\mu(S_x)=0$ as a literal identity of chain maps.
 Indeed Cayley--Hamilton holds on the finite bundle before tensoring
 with the arbitrary dg module $\mathcal A$ in that proposition;
 tensoring and then taking exact rational invariants preserve the
 identity.  In particular this assertion requires no finiteness
 assumption on $\mathcal A$ or on the whole spherical cohomology.

 Now the completed complex
 \[
    \widehat C_{\mu,v}^{\mathrm{sph}}
          =C_{\mu,v}^{\mathrm{sph}}\otimes_{Z_x}R_x
 \]
 is defined using that actual dg module structure.
 The residual roots of $P_\mu$ are the products of weight values
 in Equation~\eqref{ext:graded-slide} at $t_x$, with their tensor
 multiplicities.  Their set is Weyl invariant, so it is the same
 root set used for $Q$.
 Factor $P_\mu=P_+P_-$ over $R_x$ by Hensel's lemma, separating the
 same maximal-modulus residual roots from the others.  For a Bezout
 identity $uP_++vP_-=1$, the chain map
 \[
                  e_+=v(S_x)P_-(S_x)
 \]
 is idempotent: $e_+^2-e_+$ is a multiple of $P_\mu(S_x)=0$.
 Its image $E_{\mu,v}$ is an actual dg direct summand.
 The coefficients lie in $R_x$, so this projector commutes with
 tensoring by any finite Koszul complex over $R_x$.

 We verify that its cohomology is the transferred extreme part;
 polynomial annihilators on cohomology alone would not have
 sufficed for the preceding chain construction.  On
 $H^*(\widehat C_{\mu,v}^{\mathrm{sph}})$ both $Q(S_x)$ and
 $P_\mu(S_x)$ vanish, and the two polynomial projectors commute.
 On the intersection of the $Q_+$ part with the $P_-$ part,
 $Q_+(S_x)$ and $P_-(S_x)$ vanish.  Their residual roots have
 different moduli, so these two polynomials are relatively prime
 over $R_x$; their Bezout identity makes this intersection zero.
 The same argument interchanging $+$ and $-$ kills the other
 cross intersection.  Thus $e_+$ and the transferred
 $\varepsilon_{\mathrm{ext}}$ act identically on cohomology.
 This argument applies to arbitrary modules, without finite
 generation.  Consequently $H^*(E_{\mu,v})=N_v$, as claimed.
\end{proof}

\begin{lemma}[Top cohomology under a Koszul test]\label{ext:koszul}
 Let $(R,\mathfrak m)$ be a noetherian local ring and let $C$ be a
 complex of $R$-modules bounded above in cohomology.  Suppose that
 $N$ is its highest nonzero cohomological degree and $H^N(C)$ is a
 finite nonzero $R$-module.  If $f_1,\ldots,f_r$ generate
 $\mathfrak m$, then
 \[
 H^N\bigl(C\otimes_R K_R(f_1,\ldots,f_r)\bigr)\ne0,
 \]
 where the Koszul complex is placed in degrees $[-r,0]$.
 No regular-sequence hypothesis is required.
\end{lemma}

\begin{proof}
 The bounded Koszul complex is a complex of finite free modules, so its
 tensor product has a convergent spectral sequence
 \[
 E_2^{p,q}=H^p\bigl(K_R(f;H^q(C))\bigr)
       \ \Longrightarrow\ H^{p+q}(C\otimes_R K_R(f)).
 \]
 Its term at $(p,q)=(0,N)$ is
 $H^N(C)/\mathfrak mH^N(C)$, which is nonzero by Nakayama's lemma.
 No later differential enters this term, since its source would have
 second index greater than $N$; none leaves it, since its target
 would have positive Koszul degree.  The term survives in total
 degree $N$.
\end{proof}

\begin{proof}[Proof of Theorem~\ref{ext:discrete}]
 For semisimple rank zero, the group is a split torus.  Its unitary
 local characters have zero absolute-value exponent, so the assertion
 holds with trivial $\phi$.  Assume the result in lower semisimple rank
 and choose a violating weight $t_0$ as above.  All the preceding
 constructions are then available for arbitrarily large $m$ in
 Equation~\eqref{ext:exact-ray}.

 Use Proposition~\ref{ext:extreme-transfer} to choose $v$, and then
 Proposition~\ref{ext:spherical-projector} to obtain the dg summand
 $E_{\mu,v}$.  Let $N$ be its highest nonzero cohomological degree.
 It exists because $H^*(E_{\mu,v})=N_v$ has finite-dimensional total
 graded space, and it satisfies
 \begin{equation}\label{ext:top-lower}
       N\ge 2dm\|\lambda_0\|^2+b,
       \qquad b=2\log_q|\zeta(t_{0,y})|.
 \end{equation}
 Use the invariant algebra generators $a_1,\ldots,a_r$ fixed before
 Lemma~\ref{amp:specialization}, and put $f_j=a_j-a_j(t_x)$.
 These functions generate the maximal ideal of $Wt_x$ in $Z_x$,
 and their images generate the
 maximal ideal of $R_x$.  Apply their finite Koszul complex to the
 actual $R_x$-dg summand $E_{\mu,v}$.
 Lemma~\ref{ext:koszul} preserves a nonzero class in degree $N$.
 The chain projector of Proposition~\ref{ext:spherical-projector}
 commutes with this tensor product, so the tested summand is a
 direct summand of the tested completed spherical complex.

 Completion can be removed from the latter test.  Indeed it is
 flat, and each $f_j$ acts null-homotopically on its Koszul complex.
 The tested cohomology is therefore annihilated by the maximal
 ideal, and extension of scalars from $Z_x$ to $R_x$ changes
 none of its cohomology groups.  The natural map from the
 uncompleted spherical Koszul test to its completion is thus
 a quasi-isomorphism.  On this specialized cohomology the polynomial
 coefficients of the chain projector reduce to their residue-field
 values.  Its image is exactly the maximal-modulus group for the
 holonomy slide $S_x$, by its Hensel factorization.  Raising an
 invertible slide to the positive power $\kappa$ preserves the
 ordering of eigenvalue moduli.  Hence the surviving degree $N$
 belongs to the extreme summand to which
 Theorem~\ref{amp:amplitude} applies, giving
 \begin{equation}\label{ext:top-upper}
      N\le C+dm\max_{e\in E_{t_x}}
                                  \langle\lambda_0,h_e\rangle.
 \end{equation}

 The restriction of $t_x$ to $L^0$ is that of $t_0$.  Hence $t_x$ differs
 from $t_0$ only by a connected-central factor.  The adjoint eigenspace
 $E_{t_x}$, its compatible neutral cocharacters, and the semisimple
 exponent are consequently independent of that factor.  Because
 $t_0$ is a violation and its central exponent is zero, the strict part
 of Theorem~\ref{amp:amplitude} gives
 \[
     \delta=2\|\lambda_0\|^2-
          \max_{e\in E_{t_x}}\langle\lambda_0,h_e\rangle>0.
 \]
 The same theorem makes $C$ uniform under the possible connected-central
 changes of $t_x$.  Its hypotheses also allow the fixed auxiliary label
 and all the finitely many hyperspecial variants; replace their
 constants by their maximum.  No level, auxiliary kernel, or
 semisimple holonomy datum varies with $m$.

 Subtracting Equation~\eqref{ext:top-upper} from
 Equation~\eqref{ext:top-lower} now yields
 $dm\delta\le C-b$, impossible for arbitrarily large $m$.
 There is no violating weight.  This completes the induction and the
 proof of Theorem~\ref{ext:discrete}.
\end{proof}
